\section{A componentwise expansion with particle records}
\label{sec:cluster-expansion}

This section supplies the algebra to which the estimates of
\Cref{thm:dhm-package} will be applied.  In particular, the connected
coefficients used below are the factorial cumulants of the pasted process,
up to explicitly controlled variation errors.  They are not a factorization
assumption at an intermediate time. The expansion of incompatibilities
between autonomously evolved physical trajectory clusters follows the
viewpoint of~\cite{BGSSCluster2022}. Here the construction must also
retain the restarted dynamical caps and the observation records;
these additions are treated explicitly below.

Put
\begin{equation}\label{eq:cluster-threshold}
 M_i=\lfloor\Lambda_i/2\rfloor.
\end{equation}
We take $\eps$ small enough that $M_i\ge2$ for every $i$.  The particle cap
$\Lambda_i$ and collision-excess cap $\Gamma$ belong to the dynamics.  The
cluster-number cap $M_i$ belongs only to the expansion.  All the operators
defined below are defined at every finite root order.  The bounds $A_i$ will
be imposed only when estimating these operators.

\subsection{Autonomous clusters and a finite forest identity}

Fix one slab, temporarily suppress its index, and let $V$ be a finite
labelled particle set.  A particle state may carry previously observed
test values, or their running sum, as extra coordinates called records;
the acceptance rule uses only positions, velocities, and the collision
prehistory on this slab.  For $C\subseteq V$, write $\Psi_C$ for the
autonomous truncated evolution of the particles in $C$, including any
prescribed observation updates.  Let $\kappa_C$ be the indicator that its
accepted collision graph on the slab is connected.  Singletons have
$\kappa_C=1$; a set with more than $\Lambda$ particles has
$\kappa_C=0$.

For disjoint $C,D$, let $a(C,D)$ indicate that their autonomous histories
have an incoming cross-contact that would be accepted with their respective
prefix collision prehistories.  Put $\chi(C,D)=1-a(C,D)$.  These are
indicators, rather than numbers of contacts.  All equalities concerning
histories below are on the common full-measure set where the finite
subsystem evolutions are defined.  Taking the union of the exceptional
sets over all subsystems of a fixed finite configuration does not change
this convention.

\begin{lemma}[Autonomous partition]\label{lem:autonomous-partition}
For a finite set $V$,
\begin{equation}\label{eq:autonomous-partition}
 \sum_{\mathcal C\in\mathcal P(V)}
 \prod_{C\in\mathcal C}\kappa_C
 \prod_{\{C,D\}\subseteq\mathcal C}\chi(C,D)=1.
\end{equation}
Here $\mathcal P(V)$ denotes the partitions of $V$.  If $S\subseteq V$ is a
root set and $H$ tests the root histories, then
\begin{equation}\label{eq:autonomous-observation}
 H(\Psi_V|_S)=
 \sum_{\mathcal C\in\mathcal P(V)}
 \prod_{C\in\mathcal C}\kappa_C
 \prod_{\{C,D\}\subseteq\mathcal C}\chi(C,D)
 H(\Psi_{\mathcal C}|_S),
\end{equation}
where $\Psi_{\mathcal C}$ is the juxtaposition of the autonomous histories.
\end{lemma}

\begin{proof}
The collision components of the full truncated history give one summand
equal to one.  Its restrictions to those components are autonomous:
no accepted collision joins different components, and the rule at an
internal contact depends only on the prefix component involved there.
Conversely, suppose the factors in a summand equal one.  Juxtapose the
corresponding autonomous histories and examine their incoming contacts in
chronological order.  The internal accepted contacts obey the full rule.
Every cross-contact fails that rule by $\chi(C,D)=1$.  Thus the juxtaposed
history obeys the full truncated dynamics.  Uniqueness identifies it with
the full history and identifies its collision components with
$\mathcal C$.  This proves both assertions.  Rejected contacts are not
inserted in a prefix collision graph in this argument.
\end{proof}

The following lemma is purely finite algebra.  Its vertices will be
autonomous clusters, not particles.

\begin{lemma}[Rooted forest stopping]\label{lem:rooted-forest}
Let $W$ be a finite vertex set, let $R\subseteq W$ be nonempty, and attach
$a_e\in\{0,1\}$ to each unordered pair of vertices.  Fix $M\ge2$ and set
\begin{equation}\label{eq:ursell-block}
 u(B)=\sum_{\substack{G\text{ connected}\\V(G)=B}}
             \prod_{e\in E(G)}(-a_e),\qquad u(\{w\})=1.
\end{equation}
There is a finite signed remainder $\mathcal R$ such that
\begin{align}\label{eq:finite-rooted-identity}
 \prod_{e\subseteq W}(1-a_e)
 &=\sum_{I:\,R\subseteq I\subseteq W}
   \sum_{\substack{\mathcal B\in\mathcal P(I)\\
               B\cap R\ne\varnothing,\ |B|<M\ (B\in\mathcal B)}}
       \prod_{B\in\mathcal B}u(B)
       \prod_{e\subseteq W\setminus I}(1-a_e)+\mathcal R.
\end{align}
Its absolute value is bounded by a sum of rooted forests.  Each forest in
this bound has exactly one component with $M\le |B|\le2M-2$; its other
components have fewer than $M$ vertices.  All these components meet $R$.
Vertices outside the forest retain their full mutual compatibility
product.  The coefficient of each selected forest in this domination is
at most one.  Moreover,
\begin{equation}\label{eq:ursell-tree-bound}
 |u(B)|\le\sum_{T\text{ spanning tree on }B}\prod_{e\in E(T)}a_e.
\end{equation}
\end{lemma}

\begin{proof}
Choose a total order on pairs.  Start with the empty selected forest and
all pair factors unprocessed.  At each step choose the first unprocessed
pair joining distinct forest components, at least one of which meets
$R$.  Expand its factor as $1-a_e$.  The positive branch removes the
factor; the negative branch selects $e$ and merges the two components.
Never expand a factor internal to a current forest component.  Stop a
branch as soon as a component reaches $M$ vertices.  If there is no
eligible pair, complete the branch.

A selected edge always joins two distinct components.  Hence selected
edges form a forest, and every nontrivial selected component is rooted.
An unrooted vertex cannot join another unrooted vertex.  At a stop, both
components just merged had fewer than $M$ vertices, which gives the
stated size range and shows that there is exactly one large component.
The selected forest determines its branch: replay the deterministic
rule, choosing the negative branch precisely on its edges.  This proves
the multiplicity assertion.

For a completed branch, expand the retained internal factors.  Each full
graph thereby produced has a unique exploration branch, obtained by
replaying the same rule and selecting precisely those eligible edges
present in the graph.  Its components meeting $R$ are precisely the
reached components of the branch.  Grouping by these components gives
the products of \eqref{eq:ursell-block}; the remaining vertices retain
their compatibility product.  If all these components have size less
than $M$, this replay cannot stop prematurely.  This proves the main
term in \eqref{eq:finite-rooted-identity}.  In particular, after grouping,
priorities involving vertices outside a block do not affect its
coefficient $u(B)$.

For a stopped branch, take absolute values and drop every unused factor
except those internal to unreached vertices.  All dropped factors are
in $[0,1]$.  The branch multiplicity just proved gives the asserted
forest bound.

For completeness, the ordinary tree inequality can be seen without a
stopping rule.  Fix a total edge order on $B$ and associate to each
connected graph its unique minimum spanning tree.  For a tree $T$, let
$Q(T)$ be the set of additional edges whose order is later than every
edge on the corresponding $T$-path.  The graphs with minimum spanning
tree $T$ are exactly $T\subseteq G\subseteq T\cup Q(T)$.  Therefore
\begin{equation}\label{eq:penrose-block}
 u(B)=(-1)^{|B|-1}\sum_T
       \prod_{e\in T}a_e\prod_{e\in Q(T)}(1-a_e).
\end{equation}
Dropping the second product proves \eqref{eq:ursell-tree-bound}.
\end{proof}

\subsection{An exact one-slab operator on factorial measures}

Write $F_s$ for ordered factorial measures normalized by $\mu^s$.
States can include records on a specified subset of the $s$ labels.
One convenient convention for their unmarked definition uses symmetric
Janossy measures $p_n$:
\begin{equation}\label{eq:janossy-factorial}
 \sum_{n\ge0}\frac{p_n(E^n)}{n!}=1,\qquad
 F_s=\mu^{-s}\sum_{n\ge s}\frac{1}{(n-s)!}
                (\operatorname{pr}_{[s]})_\#p_n.
\end{equation}
Thus $p_n/n!$ is the law restricted to $N=n$.  This notation does not
impose a density assumption.

For $S=[s]$ and $V=[s+j]$, define the signed root-output kernel
\begin{align}\label{eq:retained-kernel}
 \mathcal K_{S,V}(X,\dd X'_S)
 =\sum_{\mathcal C\in\mathcal P(V)}\prod_{C\in\mathcal C}\kappa_C(X)
 \sum_{\substack{\mathcal B\in\mathcal P(\mathcal C)\\
       (\bigcup_{C\in B}C)\cap S\ne\varnothing,\ |B|<M}}
    \prod_{B\in\mathcal B}u(B;X)\,
    \delta_{\Psi_{\mathcal C}(X)|_S}(\dd X'_S).
\end{align}
The inner restrictions apply to every $B\in\mathcal B$; its cardinality
counts clusters.  The incompatibilities in $u$ are those of the fixed
autonomous histories.  The retained operator is
\begin{equation}\label{eq:retained-operator}
 (\mathcal L_sF)(\dd X'_S)
   =\sum_{j=0}^{s(M-1)\Lambda-s}\frac{\mu^j}{j!}
       \int F_{s+j}(\dd X_V)\,
                       \mathcal K_{S,V}(X,\dd X'_S).
\end{equation}
An empty summation restriction in these formulas gives zero.  There are
at most $s$ retained blocks, so the indicated bound on $j$ follows from
the component caps.

For the stopped kernel, $V$ consists only of reached particle labels;
outside labels have already been designated as spectators.  Let
$\mathfrak F^{\mathrm{stop}}_S(\mathcal C)$ be the forests spanning the
entire vertex set $\mathcal C$ such that every forest component contains
a cluster meeting $S$, exactly one component has between $M$ and $2M-2$
cluster vertices, and every other component has fewer than $M$ cluster
vertices.  In particular, an unrooted isolated cluster is not allowed.
Define $\mathcal K^{\mathrm{stop}}_{S,V}$ by replacing the inner block
sum and its product of $u(B;X)$ in \eqref{eq:retained-kernel} with
\[
 \sum_{F\in\mathfrak F^{\mathrm{stop}}_S(\mathcal C)}
                            \prod_{e\in E(F)}a_e(X).
\]
Keep $\prod_C\kappa_C$ and the same root-output map.  The mutual
compatibility factors of the spectator clusters are summed using
\eqref{eq:autonomous-partition} in the proof below; they are not omitted
from the finite stopped-branch identity.  Define $\mathcal E_sF$ as in
\eqref{eq:retained-operator}, with this kernel and the larger upper bound
\begin{equation}\label{eq:stopped-input-order}
 j\le (s+1)(M-1)\Lambda-s.
\end{equation}

\begin{lemma}[One-slab identity and variation remainder]
\label{lem:one-slab-expansion}
For a positive input point-process law, let $F'$ be its factorial
hierarchy after the truncated slab.  If its factorial measures up to the
order in \eqref{eq:stopped-input-order} are finite, then
\begin{equation}\label{eq:one-slab-exact}
 F'_s=\mathcal L_sF+R_s,\qquad |R_s|\le\mathcal E_sF.
\end{equation}
The same statements hold with prescribed records on root labels.
The positive expression $\mathcal E_sF$ is a domination of $R_s$, not an
assertion that the signed remainder has this local kernel.
\end{lemma}

\begin{proof}
Apply \eqref{eq:autonomous-observation} and
\Cref{lem:rooted-forest} inside \eqref{eq:janossy-factorial}.  For a
completed branch, fix the reached particle labels $V\supset S$, of size
$s+j$.  The cluster partitions on the unreached labels sum to one by
\eqref{eq:autonomous-partition}, pointwise before integration against
$p_n$.  No independence of the input is used.  Choosing the reached
labels gives
\begin{equation}\label{eq:fresh-label-factorial}
 \frac{1}{(n-s)!}\binom{n-s}{j}
      =\frac{1}{j!(n-s-j)!}.
\end{equation}
Summing the remaining labels gives $\mu^{s+j}F_{s+j}$; the output
normalization $\mu^{-s}$ leaves exactly $\mu^j/j!$.

For the stopped terms, retain the mutual compatibility of unreached
clusters and discard the other unused factors and the branch-selection
restrictions.  Positivity of the input and the multiplicity bound give
$|R_s(H)|\le(\mathcal E_sF)(|H|)$ for every bounded measurable $H$.
The same spectator summation and \eqref{eq:fresh-label-factorial} apply.
There are at most $s$ rooted components, only one of which has up to
$2M-2$ clusters.  Their total number of clusters is at most
$(s+1)(M-1)$, proving \eqref{eq:stopped-input-order}.

These bounds also justify the Janossy summations: the dominating
expression involves only finitely many finite factorial measures.
Finally, the spectator cancellation and compatibility factors do not
use records, and root-output maps carry them without alteration except
for the prescribed updates.  This proves the record version.
\end{proof}

For a finite signed measure $\nu$, write
$\|\nu\|_{\TV}=|\nu|(\text{whole space})$ for its total variation norm.
We will also use a deliberately crude operator bound.  If $\Delta$ is a
signed family of input measures, then, uniformly in the record choices,
\begin{equation}\label{eq:crude-operator}
 \|\mathcal L_{i,s}\Delta\|_{\TV}
 +\|\mathcal E_{i,s}|\Delta|\|_{\TV}
 \le e^{C_*A_i\Lambda_i^2\log(1/\eps)}
       \max_{q\le A_{i-1}/2}\|\Delta_q\|_{\TV},
 \qquad s\le A_i/2.
\end{equation}
Indeed every relevant $q$ is at most $(s+1)\Lambda_i^2/2$, which is less
than $A_{i-1}/2$ for small $\eps$.  There are at most $q^q$ particle
partitions and $q^q$ cluster partitions; the spanning-tree or forest
sum is at most $q^{2q}$.  Their kernels are pushforwards multiplied by
indicators.  The activity is at most $\mu^q$ and there are at most
$A_i\Lambda_i^2$ values of $q$.  Taking logarithms proves
\eqref{eq:crude-operator}.  This estimate concerns variation only; it
asserts no Gaussian weight for a general input error.

\subsection{The one-root map and its Boltzmann approximation}

Let $g_0$ be the exact normalized initial intensity from
\Cref{prop:initial}.  Define the unmarked one-root measures recursively by
\begin{equation}\label{eq:one-root-definition}
 g_i=\mathcal U_i(g_{i-1}),\qquad
 \mathcal U_i(g):=\mathcal L_{i,1}\bigl((g^{\otimes q})_{q\ge0}\bigr).
\end{equation}
The definition is finite and makes sense first as a signed measure.  It
uses the entire retained one-root block in
\eqref{eq:retained-kernel}; there is no extra cutoff depending on the
number of observations.

We spell out the comparison needed to apply the one-root part of
\Cref{thm:dhm-package}.  A particle-line molecule below includes the
ordering of contacts along each line and a C or O flag at each contact.
Write $n$ for its number of particle lines and $r$ for its cycle excess.
The signed unrestricted tree operator uses sign $(-1)^{\#\mathrm O}$,
the usual hard-sphere incoming contact measure, and the original
$\mu^{n-1}/(n-1)!$ labelled normalization.  It prescribes the displayed
contacts and ignores all other contacts.

\begin{lemma}[Modified one-root bridge]\label{lem:one-root-bridge}
The difference between $\mathcal U_i(g)$ and the signed unrestricted
one-root particle-tree sum with $n<\Lambda_i/3$ is dominated, after
replacing inputs by $|g|$, by the following positive history integrals:
\begin{enumerate}[label=(\roman*)]
 \item retained one-root molecules with $r\ge1$;
 \item retained one-root molecules with $n\ge\Lambda_i/3$;
 \item the small unrestricted trees with one additional incoming O
 contact inserted between a pair of their flight segments.
\end{enumerate}
The first two families are subsets of the DHM one-root graph class.
The third is its additional-overlap error class.  Coefficients in these
dominations have at most a fixed power of the cutoffs.  The leading
tree coefficient is exactly $(-1)^{\#\mathrm O}$, with the labelled
factorial just stated.
\end{lemma}

\begin{proof}
Apply \eqref{eq:penrose-block} to the single retained block.  Refine each
autonomous C-cluster according to its complete ordered accepted
collision history.  Its internal indicator is at most the indicator
requiring those prescribed C contacts.  Every chosen cluster-tree edge
requires the existence of a prospective accepted incoming encounter
between two isolated histories.  Each history has finitely many flight
segments: a cluster of $q$ particles has at most $q-1+\Gamma$ accepted
collisions.  A fixed pair of straight segments has at most one incoming
contact.  Partition the encounter event by the first admissible segment
pair in a fixed order.  This writes it as a sum of disjoint encounter
indicators, each with an additional factor in $[0,1]$.  Ordering the
resulting contacts along particle lines gives particle-line molecules.
All indicators not requiring the displayed encounters can be kept as
an additional weight between zero and one.

The C clusters and the selected O forest determine the connected
retained block.  Its sign is $(-1)^{\#\mathrm O}$.  Its global cycle
excess is the sum of the C-cluster excesses.  Thus all cyclic terms and
all terms with $n\ge\Lambda_i/3$ have the first two dominations.  Their
cluster counts are below $M_i<\Lambda_i$, and their C-cluster caps are
the original $\Lambda_i,\Gamma$, so they belong to the indicated DHM
class.

Consider a remaining tree molecule.  Every prescribed C component is a
tree on fewer than $\Lambda_i/3$ particles.  At every prescribed contact
its acceptance test is below both thresholds.  If there are no other
incoming encounters on this prescribed history, its autonomous
C-clusters are exactly the displayed ones, every chosen O contact is
the unique contact selected on that cluster pair, and every unused
factor in \eqref{eq:penrose-block} equals one.  The added weight is
therefore one.  Consequently a discrepancy from the unrestricted tree
is supported on a further incoming encounter.

This last assertion also covers an internal history that fails the
autonomous-cluster indicator: take the first undesignated accepted
contact.  Until that contact the autonomous and prescribed histories
coincide, so the prescribed history has the required additional incoming
encounter.  The same argument applies to a discarded unused
compatibility factor or to a nonfirst choice of a cluster-pair contact.
Inserting an O atom at that encounter changes no velocity and hence
does not change the prescribed output map.  It supplies a positive
dominating history integral.  There are at most a fixed power of $n$
segment-pair choices and order insertions, which is a fixed cutoff
power.

Conversely, every small unrestricted particle tree is admitted by the
retained block rule when it has no further encounters: its C components
have size at most $n$, their excess is zero, and there are at most
$n<M_i$ clusters.  The C flags determine those clusters and the O flags
their spanning tree.  The contact ordering determines the segment
placements.  Thus it occurs with precisely the displayed sign, without
an omitted symmetry factor.  Equivalently, keep all fresh particle
labels throughout the argument: all the refinements partition indicators
and leave $1/(n-1)!$ unchanged.  This proves the lemma.
\end{proof}

\begin{proposition}[Independent factors]\label{prop:one-root}
The measures $g_i$ have Gaussian Bol-bounded density representatives.
With the exponents $\beta_i,\theta_i$ of
\Cref{thm:dhm-package}, for sufficiently small $\eps$,
\begin{equation}\label{eq:one-root-convergence}
 \|g_i-f_{i\tau}\|_{\Bol,\beta_i}\le\eps^{\theta_i},\qquad
 \sup_{0\le i\le L}\|g_i\|_{\Bol,\beta_i}\le C.
\end{equation}
The same estimates hold for the one-root map run through any portion
$\sigma\in[0,\tau]$ of slab $i$, with target
$f_{(i-1)\tau+\sigma}$ and exponent $\beta_i$.
\end{proposition}

\begin{proof}
\Cref{prop:initial} gives
$\|g_0-f_0\|_{\Bol,\beta_0}\le C\eps$, which satisfies the initial
smallness hypothesis for $\theta_0<1$.  Suppose the assertion holds at
the preceding interface.  Apply \Cref{lem:one-root-bridge}.  The
fixed-endpoint part of \Cref{thm:dhm-package} bounds its three error
families, including their cutoff-polynomial coefficients, by
$o(\eps^{\theta_i})$ in $\Bol,\beta_i$.  That same part compares the
signed small-tree sum with $f_{i\tau}$, under the preceding
$\Bol,\beta_{i-1}$ error hypothesis, with the remaining margin in
$\eps^{\theta_i}$.  This proves the first estimate by induction; the
uniform Bol bound follows from the Bol bound for $f$.

This use is a fixed-endpoint Bol estimate, not an inference from an
integrated $L^1$ estimate.  The source representation underlying the
one-root item of \Cref{thm:dhm-package} fixes the root endpoint before
integrating the other lines.  It therefore also provides the asserted
density representatives and preserves the summable spatial cell norm.
Restricting all contact times to a partial slab gives the uniform
partial-slab version of that item.  The same bridge and induction then
apply for every $\sigma\le\tau$.
\end{proof}

\subsection{Marked retained hierarchies and exact resummation}

Fix a finite list of observation pulses.  For estimates, let $m$ be a
fixed bound on both the number of pulses and the number of labels
carrying records.  For a particle label $a$, a
record is the sum of its requested test values at pulses already passed.
Only labels whose records are required carry this extra coordinate.
Thus all record coordinates lie in one fixed bounded interval.  For a
root set $S$ and $R\subseteq S$, let $F^R_{i,S}$ be the pasted factorial
measure at interface $i$, with record coordinates on $R$.  The operators
$\mathcal L_i$ act on these measures by appending the slab observations
on those same labels.  Compatibility is independent of the records.

Define the retained hierarchy by the finite composition
\begin{equation}\label{eq:retained-hierarchy}
 \widehat F^R_i=\mathcal L_i\mathcal L_{i-1}\cdots\mathcal L_1F^R_0.
\end{equation}
For a fixed output order, each factor uses only finitely many previous
orders, so this definition needs no convergence of a series in the
number of layers or particles.  The initial hierarchy is expanded by
\Cref{prop:initial} into its exact factorial cumulant blocks.  Initial
observation pulses, if any, act by the separate one-particle update maps
on these blocks.

\begin{lemma}[Resummation with specified records]
\label{lem:marked-regrouping}
The retained hierarchy admits an exact finite expansion
\begin{equation}\label{eq:dependent-subsets}
 \widehat F^R_{i,S}
   =\sum_{H:\,R\subseteq H\subseteq S}
        E^R_{i,H}\otimes g_i^{\otimes(S\setminus H)}.
\end{equation}
The coefficients can be represented by layered diagrams with the
following rule.  Every retained block in slab $i$ contains a root from
the interface above it.  A block is retained in $E^R_{i,H}$ if it
contains at least two such roots, or it meets a dependent input label
from the interface below.  All other blocks sum to the independent
factor $g_i$ on their unique root.  Every record-bearing line continues
to time zero.  Forgetting any record commutes with
\eqref{eq:retained-hierarchy}; at the diagram level it is effected by
resumming the newly unmarked one-root histories.
\end{lemma}

\begin{proof}
At time zero expand the factorial hierarchy into exact initial cumulant
blocks.  Put in $H$ every label in a nontrivial block and every label in
$R$.  An unmarked singleton contributes $g_0$ on the complement.  A
marked singleton remains in $H$ and carries its deterministic initial
record.  This proves the initial case.

Insert the preceding-interface expansion into
\eqref{eq:retained-operator}.  Let $H_-$ be its dependent input set.
Separate the blocks that have one new root and are disjoint from $H_-$.
Their inputs are independent $g_{i-1}$ factors; summing each such block
is exactly \eqref{eq:one-root-definition}.  Since $R\subseteq H_-$, these
blocks carry no required record.  Keep all the other blocks and call
their new root set $H$.  This gives \eqref{eq:dependent-subsets} and the
stated rule.  Each required record label lies in $H_-$ at every step,
so its own particle line is continued through all interfaces.

Fresh labels assigned to different blocks are disjoint and distributed
with their multinomial coefficients.  Identity
\eqref{eq:fresh-label-factorial} therefore permits the separate sums
used here.  All sums are finite; no conditional convergence or
factorization of the actual input law is involved.

Finally, forgetting records commutes with each root-output kernel in
\eqref{eq:retained-kernel}, because its weight uses only the unmarked
histories.  It also commutes with the initial one-particle updates.
Thus it commutes with their composition.  In the regrouped expression,
the blocks freed from the requirement of carrying a record are summed
using the same one-root definition.  This proves the asserted diagram
interpretation of projection.
\end{proof}

All nontrivial initial cumulant blocks are regarded as connections when
speaking of a global component.  For estimates, each is bounded by its
initial spanning trees from \Cref{prop:initial}.  More explicitly, a
block on $k$ labels is dominated by the product of its $f_0$ factors,
$(Ck)^{Ck}$, and a sum of root trees with the edge kernel from
\eqref{eq:dhm-initial-link}:
\begin{equation}\label{eq:initial-kernel-expansion}
 \ell_\eps(x,x')=\ind_{|x-x'|\le\eps}
                    +\eps^{\upsilon}\ind_{|x-x'|>\eps},
 \qquad \upsilon=1/81.
\end{equation}
One can use this piecewise version of the kernel in the initial bound:
it follows directly from the root-proximity-component proof in
\Cref{prop:initial}.  For bounded $k$, its sharper version has kernel
$\ind_{|x-x'|\le C_{D,k}\eps}+\eps^D$ for any fixed $D$.
These are dominations of the exact initial blocks.  An artificial link
used to estimate a marked line is not an initial cumulant connection.

\begin{lemma}[Labels and connected components]\label{lem:diagram-labels}
The retained expansion factors over its global components.  If
$D^R_{i,B}$ is the sum of its diagrams connected on the designated roots
$B$, then
\begin{equation}\label{eq:component-exponential-formula}
 \widehat F^R_{i,S}
       =\sum_{\pi\in\mathcal P(S)}
                         \bigotimes_{B\in\pi}D^{R\cap B}_{i,B}.
\end{equation}
Consequently factorial-cumulant inversion of the retained hierarchy
keeps exactly its connected diagrams.  This assertion is unchanged by
the resummations into $g_i$.

At fixed root order $h=|S|$ and bounded DHM complexity $\rho\le K$,
let $n_i$ denote the number of
particle lines in slab $i$, and $j_i$ its fresh input labels.  The
labelled denominator is $\prod_i j_i!$, where
\begin{equation}\label{eq:bounded-crossing-factorials}
 n_i-j_i=O_{K,h,m,L}(1),\qquad
 \prod_i\frac1{j_i!}
      \le\frac{(1+\sum_i n_i)^{C_{K,h,m,L}}}{\prod_i n_i!}.
\end{equation}
Apart from the choices of ordinary contact pairs and their order, all
distinguished-label, gluing, initial-block, and record-placement
multiplicities are bounded by a fixed polynomial in total size times
$C^b$, where $b$ is the atom count and the exponential base $C$ does not
depend on $K$ or the moment order.  In the coarse regime the same
choices cost a power of the cutoffs with exponent $O(\rho)$, in addition
to the exponential atom factor.
\end{lemma}

\begin{proof}
First expand each $g_i$ using its finite definition all the way to time
zero, only for this algebraic argument.  Each slab block and each
initial cumulant block has a weight depending only on its labels.
The restrictions on its size and excess are componentwise.  Every
global component contains a designated root, since each slab block
meets a root in the interface above it.  On distributing $j$ fresh
labels among component sizes $j_1,\ldots,j_q$, the coefficient is
\[
 \frac1{j!}\frac{j!}{j_1!\cdots j_q!}
                 =\prod_{a=1}^q\frac1{j_a!}.
\]
All integrations and all remaining restrictions consequently factor.
Grouping the components proves
\eqref{eq:component-exponential-formula}.  Ordinary finite partition
inversion gives its cumulant assertion.  A one-root resummation
contracts a history attached to the rest through one particle line at
its top.  It cannot split or merge global components, so the same
assertion holds after those resummations.

For the counting statement, the only labels in a slab that are not
fresh are the designated roots or labels continued through an
interface.  Their number is bounded when $\rho$ and the observation
list are fixed, proving the first assertion in
\eqref{eq:bounded-crossing-factorials}; the second follows from
$n!/(n-d)!\le n^d$.  Choosing or matching boundedly many such labels
costs a fixed power of total size.  The same is true of the initial
link vertices and the segments carrying a fixed list of observations.
There are $2^b$ C/O flag choices.  Once an ordered C/O graph is fixed,
its C components determine the autonomous collision clusters and its
selected O forest determines their blocks.  There is no further
arbitrary partition of all ordinary lines.  The forest exploration has
at most one branch per selected forest.  Thus no additional factorial
or order-dependent exponential base is introduced by these operations.
At unbounded complexity the same argument bounds each distinguished
choice by a power of the maximum slab size; its exponent is controlled
by the crossing and initial-link counts.  This is the stated cutoff
power, consistent with the molecule count in
\Cref{thm:dhm-package}.
\end{proof}

\subsection{Variation errors and removal of high complexity}

The next estimates use the coarse, positive history-integral item of
\Cref{thm:dhm-package}.  There is no use of a comparison between the
true and pasted dynamics in this subsection.

\begin{lemma}[Stopped variation errors]\label{lem:stopped-variation}
Choose $L$ sufficiently large in terms of the fixed data and $\Gamma$.
For any fixed bound on the number of recorded labels, there are
$c_i>0$ such that
\begin{equation}\label{eq:retained-variation-error}
 \|F^R_{i,S}-\widehat F^R_{i,S}\|_{\TV}
                  \le e^{-c_i\Lambda_i},
 \qquad |S|\le A_i/2.
\end{equation}
The constants may decrease with $i$ and the fixed record bound.
The statement is uniform for a partial final slab.  No velocity weight
is claimed for this error.
\end{lemma}

\begin{proof}
Write $\Delta_i=F_i-\widehat F_i$.  By
\Cref{lem:one-slab-expansion},
\begin{equation}\label{eq:variation-recursion}
 \Delta_i=R_i+\mathcal L_i\Delta_{i-1},\qquad
 |R_i|\le\mathcal E_iF_{i-1}.
\end{equation}
The initial difference is zero.  In the positive stopped bound substitute
$F_{i-1}=\widehat F_{i-1}+\Delta_{i-1}$, take absolute values, and
expand the retained part as in \Cref{lem:marked-regrouping}.  Use the
initial tree domination and \Cref{prop:one-root}.  Refining autonomous
histories into C graphs and selecting one incoming contact for every
incompatibility gives exactly the positive C/O histories of the coarse
package, with their previous-interface continuations.  The only
top-slab difference is a stopped component with between $M_i$ and
$2M_i-2<\Lambda_i$ clusters.  Hence the same coarse family contains it,
and its atom count is at least $M_i-1$.

Here are the bookkeeping losses in this application.  The independent
inputs from the lower-layer expansion have already been resummed into
$g_{i-1}$ (and earlier $g$'s).  We do not resum any component of the
positive stopped forest in the top slab.  Every top component that
fails the continuation condition is covered by the retained
root-deficit bound in \Cref{lem:dhm-dummies}(i), at total cost
$\eps^{-C_*|S|}$.  To satisfy the bottom continuation condition on a
required record line that has no physical initial link, use
\Cref{lem:dhm-dummies}.  Equivalently, add an independent, unit-mass,
positive Gaussian Bol input on a straight line continued from zero to
the top, and connect it to that line by
\eqref{eq:initial-kernel-expansion}.  Since this kernel is at least
$\eps^\upsilon$, insertion costs at most $\eps^{-\upsilon}$; the new
activity and top-root normalization cancel.  Only boundedly many such
lines are needed.  The $A_i/2$ slack accommodates them.  They are used
only in this upper bound.

The remaining label and initial-block factors are the cutoff powers in
\Cref{lem:diagram-labels}.  The coarse bound therefore gives, after
summing complexities and atoms,
\begin{equation}\label{eq:stopped-coarse-estimate}
 \|\mathcal E_i\widehat F_{i-1}\|_{\TV}
  \le \eps^{-C_*|S|-C_{m}}
          \sum_{b\ge M_i-1}(C\tau^{c_0})^b
          \sum_{\rho\ge0}|\log\eps|^{C_*\rho}\eps^{c_*\rho}
  \le e^{-c\Lambda_i}.
\end{equation}
In this display $c_0,c_*>0$ are the fixed coarse exponents; absorbed
constants can depend on $L$ and the fixed observation list.  Choose
$L$ so that $C\tau^{c_0}<1$.  For small $\eps$ the complexity series is
geometric.  Finally $\Lambda_i=A_i^{30}$ and $A_i\ge|\log\eps|$ imply
$A_i\log(1/\eps)=o(\Lambda_i)$, so the displayed top losses are
absorbed in the exponential size gain.

All input orders are at most $(|S|+1)\Lambda_i^2/2<A_{i-1}/2$.
The terms involving $\Delta_{i-1}$, both in the stopped bound and in
\eqref{eq:variation-recursion}, are bounded by
\eqref{eq:crude-operator}.  The induction hypothesis and
\[
 A_i\Lambda_i^2\log(1/\eps)=o(\Lambda_{i-1}),
 \qquad A_{i-1}=\Lambda_i^{30},
\]
absorb these terms into an error smaller than
$e^{-c\Lambda_i}$.  This proves the induction.  Restricting all times
in the last slab to a subinterval preserves every upper bound just
used, proving the partial-slab assertion.
\end{proof}

\begin{theorem}[Marked component expansion]\label{thm:marked-expansion}
Choose $\Gamma$ and then $L$ as above.  These choices are independent of
the moment order and the fixed observation list.  On this single pasted
dynamics the following statements hold.
\begin{enumerate}[label=(\roman*)]
 \item The factorial hierarchy with any fixed finite list of records
 has the finite componentwise expansion
 \eqref{eq:retained-hierarchy}, with the variation error
 \eqref{eq:retained-variation-error}.  Its independent inputs are the
 $g_i$ in \Cref{prop:one-root}; its label weights, continuations, and
 record projections are those of
 \Cref{lem:marked-regrouping,lem:diagram-labels}.
 \item At each fixed final order $h$, the actual normalized marked
 factorial cumulant $c_h^\eps$ is the sum of retained globally connected
 diagrams, plus an error of total variation at most
 $C_h e^{-c\Lambda_L}$.  A nontrivial initial cumulant block is part of
 this connectivity.  Dummy links are not.
 \item For any fixed accuracy power $P$, one can choose a constant
 $K$ after the observation list and $h$ such that all remaining
 retained diagrams with $\rho>K$ have summed absolute variation at
 most $\eps^P$.  At $\rho\le K$ the original fresh-label factorials
 satisfy \eqref{eq:bounded-crossing-factorials}.  Neither increasing
 $K$ nor increasing a fixed moment order changes the exponential base
 per ordinary atom in the preceding estimates.
\end{enumerate}
\end{theorem}

\begin{proof}
The first assertion combines the preceding propositions and lemmas.
For the second, apply factorial-cumulant inversion to the actual and
retained hierarchies.  Formula
\eqref{eq:component-exponential-formula} gives precisely the connected
sum for the latter.  Every term in the difference of the two finite
partition polynomials has a factor $\Delta_{L,B}$.  The actual
normalized factorial measures have mass at most one: initially this
follows from the insertion bound in \Cref{prop:initial}, and the
dynamics preserves $N$.  The retained measures therefore have
variation at most $1+e^{-c\Lambda_L}$ at fixed orders by
\Cref{lem:stopped-variation}.  This proves the stated cumulant error.

For the third assertion use the coarse bound on each resummed core,
with the same fixed record and top-root losses used in
\eqref{eq:stopped-coarse-estimate}.  Independent top singleton factors
are bounded by \Cref{prop:one-root}.  For sufficiently small $\eps$,
$|\log\eps|^{C_*\rho}\eps^{c_*\rho}\le\eps^{c_*\rho/2}$.
After summing the geometric atom factor, the sum with $\rho>K$ is at
most
\[
 C_{h,m}\eps^{-C_*h-C_m}\eps^{c_*K/2}.
\]
Increase $K$ to make this at most $\eps^P$.  Adding the boundedly many
dummy lines changes the complexity by a fixed amount and is absorbed
by the same choice.  Remove these artificial lines before applying any
low-complexity asymptotics.

All root-order and observation choices in this reasoning occur in the
cutoff powers of the complexity estimate or in the fixed polynomial
factors of \eqref{eq:bounded-crossing-factorials}.  The exponential
atom bases come from the fixed graph-counting and Bol bounds, with
$\Gamma$ fixed first.  Thus the required choice of $L$ is independent
of each subsequently fixed order.  This completes the proof.
\end{proof}

\subsection{Terminal witnesses without observed roots}

The comparison with true dynamics in \Cref{sec:cutoff} requires a
positive version of the expansion with a connected final collision
witness instead of observed outputs. We record that consequence now;
like the marked expansion, it requires no bound on the cutoff probability.

\begin{corollary}[Terminal-witness expansion]\label{cor:terminal-witness}
In slab $i$, let a terminal witness be a connected true C history on
$n\le3\Lambda_i$ incoming labels, of excess at most $2\Gamma$, stopped
at a contact time in that slab.  It may carry arbitrary further
nonnegative realizability restrictions.  Sum its indicator with the
full activity normalization
\begin{equation}\label{eq:terminal-witness-functional}
 \sum_n\frac{\mu^n}{n!}
       \int F_{i-1,n}(\dd Z_n)\,
                       \ind_{\text{terminal witness}}(Z_n).
\end{equation}
For the terminal classes in \Cref{lem:dhm-dummies}(ii), this
expression has the preceding retained
layer expansion, with independent $g$ inputs and initial blocks.
All lower C/O components attach through particle lines to the
terminal witness.  The preceding variation error is smaller than every
fixed power of $\eps$.  After the artificial-top-root normalization
cost $\mu$, the coarse part with $\rho>K$ can also be made smaller than
any assigned power by choosing $K$ sufficiently large.  The remaining
low-complexity witness integrals retain their realizability restrictions.
\end{corollary}

\begin{proof}
Here $n\le3\Lambda_i<A_{i-1}/2$ for small $\eps$, so substitute the
preceding-interface expansion from
\Cref{thm:marked-expansion} in \eqref{eq:terminal-witness-functional}.
Every lower slab component contains a root leading to the slab above,
and the final set of those lines enters the connected witness.  This
proves the attachment assertion before discarding any initial link.
The contribution of the preceding variation error is bounded by
\[
 \sum_{n\le3\Lambda_i}\frac{\mu^n}{n!}
       e^{-c\Lambda_{i-1}}
 \le e^{C\Lambda_i\log(1/\eps)-c\Lambda_{i-1}},
\]
which has the claimed smallness.  Arbitrary witness restrictions are
indicators and are kept in the positive integrals.  Designating an
artificial top root places the remaining diagrams in the stated
terminal-error class when there is a dependent continuation below;
removing its normalization costs $\mu$.  If there is no such
continuation, the diagram has only its terminal slab with independent
inputs, and its complexity is at most $1+2\Gamma$.  The same geometric
complexity summation as in the proof of
\Cref{thm:marked-expansion}, with this additional fixed power, proves
the last assertion.  For $i=1$ there is no preceding variation error;
the identical argument starts directly from the initial cumulant
expansion.  Nothing here compares true and pasted evolutions.
\end{proof}
